\section{Related Work}

\subsection{Communication Accommodation Theory}

Communication Accommodation Theory (CAT), introduced by \citet{giles1973accommodation}, posits that individuals adjust their communication style to signal social identity and regulate social distance. Convergence---adapting toward an interlocutor's style---signals affiliation and builds rapport, while divergence maintains distinctiveness.

\citet{niederhoffer2002linguistic} operationalized accommodation through Linguistic Style Matching (LSM), finding correlations around $r\approx0.3$ between style matching and relationship quality in human-human dyads. This work established that accommodation is measurable through linguistic features and predicts interaction outcomes.

However, human-human accommodation research assumes symmetric agency: both parties can choose to accommodate. In human-AI interaction, this symmetry breaks. AI systems accommodate through training objectives, not social motivation; humans may not perceive AI as warranting reciprocal accommodation.

\subsection{Human-AI Linguistic Adaptation}

Recent work confirms that accommodation occurs in human-AI dialogue. \citet{chen2026bidirectional} analyzed 1,319 GPT-4o conversations from WildChat, demonstrating bidirectional linguistic accommodation with distinct temporal dynamics: model accommodation is front-loaded (strongest in early turns), while user convergence develops gradually.

This finding validates that accommodation exists in human-AI interaction, but leaves the \emph{outcome} relationship untested. Our work extends this line by testing whether accommodation magnitude predicts conversation continuation.

\subsection{Chatbot Engagement and the Uncanny Valley}

\citet{ciechanowski2019uncanny} documented uncanny valley effects in chatbot interaction: users respond negatively to AI behaviors that approximate but don't fully achieve human-like communication. This suggests that excessive accommodation---making AI ``too human''---may backfire.

Our finding of an inverted-U accommodation-engagement relationship aligns with this framework. Moderate accommodation signals attentiveness without triggering artificiality concerns; excessive accommodation may activate uncanny valley responses.
